\documentclass[10pt,a4paper]{amsart}
\usepackage[margin=19mm]{geometry}
\usepackage{amsmath,amssymb,mathtools}
\usepackage[scr=rsfs,cal=euler]{mathalfa}
\usepackage{xfrac}
\usepackage[unicode,hidelinks,bookmarks=false]{hyperref}
\newcommand{\ie}{i.e.\ }
\newcommand{\cf}{cf.\ }
\newcommand{\resp}{respectively\ }
\newcommand{\Subsection}{\subsection}
\providecommand{\nobreakdash}{\nobreak\hbox{-}}
\setlength{\emergencystretch}{2em}
\multlinegap0pt
\usepackage{amssymb}
\usepackage{mathtools}
\usepackage{xcolor}
\newtheorem{lemma}{Lemma}
\newcommand{\trans}{^{\mathrm{T}}}
\title{Walk Matrix–Based Upper Bounds on Generalized Cospectral Mates}
\author{Muhammad Raza, Obaid Ullah Ahmad, Mudassir Shabbir, Waseem Abbas}
\date{Source: arXiv:2601.07373v3 (CC BY 4.0)}
\begin{document}
\maketitle

\noindent\emph{Source and licence.} Walk Matrix–Based Upper Bounds on Generalized Cospectral Mates. Version arXiv:2601.07373v3 (CC BY 4.0), distributed by arXiv under the Creative Commons Attribution 4.0 International licence, http://creativecommons.org/licenses/by/4.0/, as recorded in the arXiv licence metadata for this exact version and shown on the abstract page https://arxiv.org/abs/2601.07373v3. Full licence notice: \texttt{licenses/arxiv-2601.07373-CC-BY-4.0.txt}.\par\smallskip

\noindent\textit{Editorial scope.} Counted unit: the article's lemma lemma:p-out (source0.tex lines 314--320). The article's complete original proof of it is delivered with it (source0.tex lines 322--357, one \texttt{proof} environment of 36 lines). No mathematics is written, repaired or re-derived: the statement and the proof are the article's own lines, in the article's own notation, and no retained line is substituted or re-flowed.  Retained uncounted as the source-local context the proof needs: lines 219--227; lines 243--251; lines 265--274.  Nothing was omitted from any retained span, so the package carries no omission record.  No retained line contains a citation, so the wrapper carries no bibliography and no bibliography entry.  No retained line contains a figure environment, an image-inclusion command or drawing code, so no image is delivered and the package declares no asset dependency.  Editorial formatting only, in addition: an AMS \texttt{amsart} preamble that restates the article's own declarations and macros used by the retained lines, namely the environments \texttt{lemma} on separate counters, together with the standard packages those lines need.  The retained environments print in this document as Lemma 1 in order of appearance, whereas the article prints the same items with the numbering of its own full document.  The complete native source of the article as distributed in the arXiv e-print for this exact version is bundled unchanged as \texttt{sources/192495/source0.tex}.
\begin{lemma}
\label{lemma:q-props}
Let $Q \in O_n(\mathbb{Q})$, $\bar{Q} = \ell(Q)Q$ and $v_1, v_2, \dots, v_n$ be the column vectors of $\bar{Q}$. Let $p$ be a prime and $k$ be a positive integer such that $p^k \parallel \ell(Q)$. Then:
\begin{enumerate}
\item $v_i\trans v_j \equiv 0 \pmod{p^{2k}}$ for all $i, j$;
\item $p^{2k} \parallel v_i\trans v_i$ for all $i$;
\item $p^{k} \parallel v_i\trans e$ for all $i$.
\end{enumerate} 
\end{lemma}

\begin{lemma}
\label{lemma:uv}
Let $u$ and $v$ be $n$-dimensional integer vectors, and let $p$ be an odd prime and $k$ be a positive integer. If the following conditions hold:
\begin{enumerate}
\item $u \equiv \alpha v \pmod{p^k}$ where $\alpha$ is an integer not divisible by $p$,
\item $u\trans u \equiv v\trans v \equiv 0 \pmod{p^{2k}}$,
\end{enumerate}
then $u\trans v \equiv 0 \pmod{p^{2k}}$.
\end{lemma}

\begin{lemma}
\label{lemma:uv-2}
Let $u$ and $v$ be $n$-dimensional integer vectors, and let $k$ be a positive integer. If the following conditions hold:
\begin{enumerate}
\item $u \equiv \alpha v \pmod{2^k}$ where $\alpha$ is an odd integer,
\item $\operatorname{ord}_2(u\trans u) = \operatorname{ord}_2(v\trans v) = 2k$,
\item $\operatorname{ord}_2(u\trans e) = \operatorname{ord}_2(v\trans e) = k$,
\end{enumerate}
then $u\trans v \equiv 0 \pmod{2^{2k}}$.
\end{lemma}

\begin{lemma}
\label{lemma:p-out}
Let $Q_1, Q_2 \in O_n(\mathbb{Q})$ and let $p$ be a prime. Suppose that 
$\operatorname{ord}_p(\ell(Q_i)) = k \geq 1$ for $i=1,2$, and that there exists a vector $w \not\equiv 0 \pmod{p}$ such that every column 
of $\ell(Q_i)Q_i$ is a scalar multiple of $w$ 
over $\mathbb{Z}/p^k\mathbb{Z}$. Then $p \nmid \ell(Q_1\trans Q_2)$.
\end{lemma}

\begin{proof}
Let $\bar{Q}_i=\ell(Q_i)Q_i$ and let $R = \bar{Q}_1\trans \bar{Q}_2$. Since $\bar{Q}_1$ and $\bar{Q}_2$ are integral, $R$ is an integral matrix with entries $r_{ij} = u_i\trans v_j$, where $u_i$ and $v_j$ denote the $i^\text{th}$ and $j^\text{th}$ columns of $\bar{Q}_1$ and $\bar{Q}_2$, respectively. We show that $r_{ij} \equiv 0 \pmod{p^{2k}}$ for all $i, j$. By the definition of level, there exist column vectors $u'$ of $\bar{Q}_1$ and $v'$ of $\bar{Q}_2$ such that $u' \not\equiv 0 \pmod{p}$ and $v' \not\equiv 0 \pmod{p}$. Since both $u'$ and $v'$ are also scalar multiples of $w$ over $\mathbb{Z}/p^k\mathbb{Z}$, there exists an integer $\alpha \not\equiv 0\pmod{p}$ such that:
\begin{equation}
\label{eq:u'-v'}
u' \equiv \alpha v' \pmod{p^{k}},
\end{equation}
Furthermore, since every column vector $u_i$ and $v_j$ is a scalar multiple of $w$ over $\mathbb{Z}/p^k\mathbb{Z}$, we can express them as:
\begin{equation}
\label{eq:ui}
u_i \equiv c_i u' + p^{k}x_i \pmod{p^{2k}},
\end{equation}
\begin{equation}
\label{eq:vi}
v_j \equiv d_j v' + p^{k}y_j \pmod{p^{2k}},
\end{equation}
where $c_i, d_j$ are integers and $x_i, y_j$ are integer vectors.

Since $Q_1$ and $Q_2$ satisfy the conditions of Lemma~\ref{lemma:q-props}, it follows that $u_i\trans u' \equiv 0 \pmod{p^{2k}}$ and $v_j\trans v' \equiv 0\pmod{p^{2k}}$ for each column $u_i$ and $v_j$. Moreover, $\operatorname{ord}_p({u'}\trans u') = \operatorname{ord}_p({v'}\trans v') = 2k$ and $\operatorname{ord}_p({u'}\trans e) = \operatorname{ord}_p({v'}\trans e) = k$. Therefore, we can apply Lemma~\ref{lemma:uv} (for odd primes) or Lemma~\ref{lemma:uv-2} (for $p=2$) to establish that ${u'}\trans v' \equiv 0 \pmod{p^{2k}}$.

Multiplying Eq.~\eqref{eq:ui} by ${u'}\trans$ yields:
\begin{align*}
{u'}\trans u_i &\equiv {u'}\trans (c_i u' + p^{k}x_i) \pmod{p^{2k}} \\
&\equiv c_i {u'}\trans u' + p^{k}{u'}\trans x_i \pmod{p^{2k}}
\end{align*}
As ${u'}\trans u_i \equiv 0 \pmod{p^{2k}}$ and $c_i {u'}\trans u' \equiv 0 \pmod{p^{2k}}$, it follows that $p^{k} {u'}\trans x_i \equiv 0 \pmod{p^{2k}}$, which implies ${u'}\trans x_i \equiv 0 \pmod{p^{k}}$. Using Eq.~\eqref{eq:u'-v'}, we get ${u'}\trans x_i \equiv \alpha {v'}\trans x_i \pmod{p^{k}}$. As $\alpha \not\equiv 0 \pmod{p}$, we conclude that ${v'}\trans x_i \equiv 0 \pmod{p^{k}}$. Using a similar approach, we can prove that ${u'}\trans y_j \equiv 0 \pmod{p^{k}}$ for each $y_j$.

Finally, for any columns $u_i$ and $v_j$, we have:
\begin{align*}
u_i\trans v_j &\equiv (c_i u' + p^{k} x_i)\trans (d_j v' + p^{k} y_j) \\
&\equiv c_i d_j {u'}\trans v' + c_i p^{k} {u'}\trans y_j + d_j p^{k} x_i\trans v' + p^{2k} x_i\trans y_j \pmod{p^{2k}}.
\end{align*}
We established that ${u'}\trans y_j \equiv 0 \pmod{p^{k}}$ and ${v'}\trans x_i \equiv 0 \pmod{p^{k}}$. Since ${u'}\trans v' \equiv 0 \pmod{p^{2k}}$, every term on the right-hand side vanishes modulo $p^{2k}$. Thus, we conclude that $u_i\trans v_j \equiv 0 \pmod{p^{2k}}$ for all $i, j$, which implies $R \equiv 0 \pmod{p^{2k}}$. Substituting the definition of $R$, we have
$$\ell(Q_1)\ell(Q_2) Q_1\trans Q_2 \equiv 0 \pmod{p^{2k}}.$$
This implies that  $\frac{\ell(Q_1)\ell(Q_2)}{p^{2k}} Q_1\trans Q_2$ is an integral matrix, and hence $\ell(Q_1\trans Q_2) \mid \frac{\ell(Q_1)\ell(Q_2)}{p^{2k}}$. Since $p^k \parallel \ell(Q_1)$ and $p^k \parallel \ell(Q_2)$, it follows that $p^{2k} \parallel \ell(Q_1)\ell(Q_2)$. Consequently, $p \nmid \frac{\ell(Q_1)\ell(Q_2)}{p^{2k}}$, which implies $p \nmid \ell(Q_1\trans Q_2)$, completing the proof.

\end{proof}

\end{document}
